%% IACR ePrint submission
%% To compile with the official IACR template, replace this line with:
%%   \documentclass{iacrtrans}
%% and download iacrtrans.cls from https://github.com/IACR/latex
%% This file compiles with standard LaTeX as-is.
\documentclass[11pt,a4paper]{article}

% ─── Packages ──────────────────────────────────────────────────────────────
\usepackage[utf8]{inputenc}
\usepackage[T1]{fontenc}
\usepackage{booktabs}
\usepackage{array}
\usepackage{listings}
\usepackage{xcolor}
\usepackage{hyperref}
\usepackage{amsmath}
\usepackage{amssymb}

\usepackage{graphicx}
\usepackage{mdframed}

% ─── Listings style (protocol / pseudocode blocks) ─────────────────────────
\lstdefinestyle{protocol}{
  basicstyle=\ttfamily\small,
  breaklines=true,
  frame=single,
  framerule=0.4pt,
  rulecolor=\color{gray!60},
  backgroundcolor=\color{gray!6},
  xleftmargin=1em,
  xrightmargin=0.5em,
  aboveskip=0.8em,
  belowskip=0.8em,
  columns=fullflexible,
  keepspaces=true,
  literate={→}{{$\rightarrow$}}1
           {−}{{$-$}}1
           {≤}{{$\leq$}}1
}

% ─── Architectural note box ────────────────────────────────────────────────
\newmdenv[
  linecolor=blue!40,
  backgroundcolor=blue!4,
  linewidth=1pt,
  leftmargin=0pt,
  rightmargin=0pt,
  innerleftmargin=8pt,
  innerrightmargin=8pt,
  innertopmargin=6pt,
  innerbottommargin=6pt,
  skipabove=\medskipamount,
  skipbelow=\medskipamount
]{archbox}

% ─── Metadata ──────────────────────────────────────────────────────────────
\title{EIDA: Ephemeral Identity Distribution Architecture\\[0.4em]\large\textit{Whitepaper v1.4}}

\author{Güven Acar}
\date{May 2026}

\begin{document}

\maketitle

% ═══════════════════════════════════════════════════════════════════════════
\begin{abstract}
Classical Public Key Infrastructure (PKI) provides static identity
guarantees---it proves that a public key belongs to an entity, but cannot
guarantee that a given transaction originates from that entity at that
moment.  This gap enables replay attacks, key-reuse vulnerabilities, and
long-term compromise cascades.

EIDA (Ephemeral Identity Distribution Architecture) addresses this gap by
introducing \emph{identity-level forward secrecy}.  The core primitive,
\textbf{ePriPub}, is a single-use ephemeral key pair derived per transaction
from the user's permanent identity inside a hardware-backed isolation zone
(TPM~2.0 / TEE).  Once used, the ePriPub components are destroyed and the
Certificate Authority (CA) issues a fresh single-use certificate---no
persistent state is required.

EIDA shifts the trust model from centralised authority verification to
cryptographic transaction uniqueness, distributing security to the user
without requiring trusted intermediaries.
\end{abstract}

\noindent\textbf{Keywords:} ephemeral keys, forward secrecy, PKI, hardware security, replay protection, identity architecture

% ═══════════════════════════════════════════════════════════════════════════
\section{Problem}

Today, companies responsible for securing user identities are under constant
attack.  A single breach of a centralised security system can expose the
personal data of millions of users.  This is an inevitable consequence of the
current architecture: security is concentrated, and concentration creates
high-value targets.

EIDA proposes distributing this responsibility to user devices.  In a
decentralised model, an attacker who compromises a single node gains access to
only one person's data---not millions.  Large technology companies are no
longer the primary target.  The attack surface shrinks from a single critical
system to millions of individual devices.

Classical PKI reinforces this centralisation problem.  It guarantees that a
public key belongs to an entity, but cannot verify that a specific transaction
originates from that entity at that moment---enabling replay attacks, key
reuse, and long-term compromise cascades.

% ═══════════════════════════════════════════════════════════════════════════
\section{Architecture}

EIDA consists of three layers.

\paragraph{1.~User Isolation Zone.}
The user's permanent private key (\texttt{uPriv}) resides exclusively within
a hardware-backed isolation zone (TPM~2.0 / TEE).  It never leaves this
environment under any circumstance.  All cryptographic operations involving
\texttt{uPriv} are performed inside the zone.

\paragraph{2.~ePriPub Derivation Layer.}
For each transaction, a single-use ephemeral key pair is derived inside the
isolation zone:
\begin{lstlisting}[style=protocol]
uPriv + nonce + timestamp → HKDF → seed → Ed25519 → (ePriv, ePub)
\end{lstlisting}
The nonce is randomly generated per transaction and combined with a
timestamp, guaranteeing that no two derivations produce the same key pair.
Transaction uniqueness is ensured by the \texttt{(ePub, nonce, timestamp)}
triple---no external token coordinator is required.

\paragraph{3.~Stateless CA Layer.}
The CA operates as a reactive signer---it does not scan, store, or track
previous requests.  Each incoming request is treated as a self-contained,
atomic unit:
\begin{lstlisting}[style=protocol]
Receive:  (ePub + nonce + timestamp + signature)
Verify:   signature valid?        → yes
          timestamp within TTL?   → yes
Generate: new (sPriv, sPub) key pair
Issue:    sCert = {sPub + expiry + metadata}, signed with CA key
Send:     (sCert + signed document) to user isolation zone
Forget:   sPriv destroyed, no state retained
\end{lstlisting}
The CA does not maintain a burned list or any persistent record of processed
requests.  Replay protection is provided by the short TTL window and the
inherent single-use nature of \texttt{sCert}.  This stateless design enables
unlimited horizontal scaling---each CA node operates independently without
database synchronisation.

% ═══════════════════════════════════════════════════════════════════════════
\section{Scalability Model: Stateless Verification and the End of Revocation
  Lists}
\label{sec:scalability}

\subsection{The Legacy Bottleneck}

Traditional PKI and identity systems suffer from what can be termed the
\emph{State Exhaustion Problem}.  To ensure that a token or key has not been
reused or compromised, a Certificate Authority must maintain large,
continuously growing Certificate Revocation Lists (CRLs) or respond to
real-time status queries (OCSP\@).  This creates a single point of latency
and a database overhead that scales poorly as user populations grow into the
hundreds of millions.

\subsection{The EIDA Model: Zero-State Scalability}

EIDA eliminates the need for a burned list or any centralised database lookup
during verification.  Rather than scaling the CA's state-management capacity,
EIDA moves the burden of security from the CA to \emph{time} and
\emph{hardware}.

\paragraph{1.~Hardware-Enforced Ephemerality.}
The ePriPub pair is generated inside a hardware-backed isolation zone
(TPM/TEE) and is strictly bound to a single transaction context.  The key
pair ceases to exist on the client side immediately after signing---not by
policy, but by design.  There is nothing to revoke because there is nothing to
persist.

\paragraph{2.~Temporal Auto-Invalidation.}
Every EIDA token is cryptographically bound to a high-precision timestamp.
The CA does not revoke keys; it simply rejects any request where
\[
  T_{\text{current}} - T_{\text{token}} > \mathit{TTL}.
\]
Invalidation is a mathematical consequence of time, not a database operation.

\paragraph{3.~$O(1)$ Verification Complexity.}
The CA performs a purely mathematical verification: it checks the signature
and the clock.  No database is queried, no list is consulted.  Verification
complexity is $O(1)$---the operation takes the same time whether the system
serves $10^{2}$ or $10^{11}$ users.

\begin{archbox}
\textbf{Architectural Note.}~~In EIDA, security is not a record maintained in
a database---it is a property of the mathematics and the hardware.  The CA
functions as a \emph{stateless gatekeeper}, making the system structurally
immune to the scalability bottlenecks that constrain modern PKI deployments.
\end{archbox}

% ═══════════════════════════════════════════════════════════════════════════
\section{ePriPub --- Ephemeral Primary Public Key Primitive}

ePriPub is the core primitive of EIDA.  It is not a protocol---it is a
single-use cryptographic identity unit that carries both signing and
verification capability for exactly one transaction.

\subsection{Definition}

ePriPub is an ephemeral key pair $(\mathit{ePriv},\,\mathit{ePub})$ where:
\begin{itemize}
  \item $\mathit{ePriv}$ --- ephemeral private component, used once to sign
        the consent token, destroyed immediately after.
  \item $\mathit{ePub}$ --- ephemeral public component, sent to the CA for
        signature verification, discarded after use.
\end{itemize}
Neither component is stored anywhere after the transaction completes.

\subsection{Derivation}

\begin{lstlisting}[style=protocol]
Inputs:
  uPriv     — permanent user private key (never leaves isolation zone)
  nonce     — 32-byte cryptographically random value (per transaction)
  timestamp — Unix timestamp in milliseconds (per transaction)

Process:
  seed           = HKDF(uPriv || nonce || timestamp)   [RFC 5869]
  (ePriv, ePub)  = Ed25519_derive(seed)                [RFC 8032]

Output:
  (ePriv, ePub)  — valid for this transaction only
\end{lstlisting}

\subsection{Lifecycle}

\begin{lstlisting}[style=protocol]
 1. DERIVE   — (ePriv, ePub) generated inside isolation zone
 2. SIGN     — consent token (ePub + nonce + timestamp) signed with ePriv
 3. TRANSMIT — (ePub + nonce + timestamp + signature) sent to CA
 4. VERIFY   — CA verifies signature and timestamp TTL
 5. DESTROY  — ePriv destroyed inside isolation zone
 6. GENERATE — CA generates new session key pair (sPriv, sPub)
 7. SIGN-DOC — CA signs document with sPriv, issues sCert
 8. RESPOND  — (sCert + signed document) sent to user isolation zone
 9. FORWARD  — isolation zone forwards to requesting institution
10. CLEAN    — sPriv destroyed by CA
\end{lstlisting}

\subsection{Security Properties}

\begin{itemize}
  \item \textbf{Single-use} --- \texttt{sCert} expires after TTL; replay
        attacks are blocked by the timestamp.
  \item \textbf{Identity-bound} --- the derivation chain links ePriPub to
        \texttt{uPriv} without exposing it.
  \item \textbf{Forward secrecy} --- compromise of any ePriPub reveals
        nothing about \texttt{uPriv} or other transactions.
  \item \textbf{Stateless CA} --- the CA retains no state between requests,
        enabling unlimited horizontal scaling.
\end{itemize}

\subsection{ePriPub as User Consent Proof}

In classical e-signature models, the user directly signs the document with
their private key.  EIDA separates this into two distinct responsibilities:

\begin{lstlisting}[style=protocol]
User  → proves consent    (ePriPub)
CA    → signs document    (sPriv)
\end{lstlisting}

The CA generates a session key pair $(\mathit{sPriv},\,\mathit{sPub})$ once
per transaction and handles the actual document signature.  However, the CA
will only proceed with signing if a valid, unused ePriPub is present in the
request.  Therefore ePriPub serves as a \textbf{cryptographic consent token}:

\begin{lstlisting}[style=protocol]
Valid ePriPub received → user consented  → CA signs and issues sCert
No ePriPub            → no consent proof → CA rejects
Expired timestamp     → TTL window past  → CA rejects
\end{lstlisting}

\paragraph{Non-repudiation guarantee.}
The user cannot later deny having initiated the transaction because:
\begin{itemize}
  \item ePriPub is derived from \texttt{uPriv}---only the user could have
        produced it.
  \item ePriPub is single-use---it cannot have been captured and replayed.
  \item The \texttt{(ePub, nonce, timestamp)} combination is unique---it
        cannot be reused in another context.
\end{itemize}
The CA cannot forge user consent because it cannot produce a valid ePriPub
without access to \texttt{uPriv}, and \texttt{uPriv} never leaves the user's
isolation zone.

% ═══════════════════════════════════════════════════════════════════════════
\section{Security Analysis}

\subsection{Threat Model}

EIDA considers the following adversaries:
\begin{itemize}
  \item \textbf{Network attacker} --- intercepts transmitted data.
  \item \textbf{CA attacker} --- compromises the certificate authority.
  \item \textbf{Device attacker} --- gains physical or logical access to the
        user's device.
  \item \textbf{Replay attacker} --- attempts to reuse captured transactions.
\end{itemize}

\subsection{Attack Scenarios}

\paragraph{Network interception.}
\begin{lstlisting}[style=protocol]
Attacker captures: (ePub + nonce + timestamp + signature)
Replay attempt   → CA rejects: timestamp TTL already expired
uPriv            → never transmitted → not exposed
Result: attack fails
\end{lstlisting}

\paragraph{CA compromise.}
\begin{lstlisting}[style=protocol]
Attacker gains access to CA
CA stores: no persistent user data, no burned list
sPriv    → destroyed after each transaction → not recoverable
uPriv    → never reaches CA → not exposed
Result: attacker finds nothing of value
\end{lstlisting}

\paragraph{Device compromise.}
\begin{lstlisting}[style=protocol]
Attacker gains access to user device
uPriv            → inside TPM/TEE → cannot be extracted
Past ePriPub     → destroyed → not recoverable
Result: attacker limited to future transactions
        only while physically holding the device
\end{lstlisting}

\paragraph{Replay attack.}
\begin{lstlisting}[style=protocol]
Attacker captures valid (ePub + nonce + timestamp + signature)
CA checks timestamp → TTL window expired → rejected
Even within TTL: sCert already delivered and consumed by institution
Institution validates sCert expiry → rejected
Result: attack fails
\end{lstlisting}

\subsection{Limitations}

\begin{itemize}
  \item EIDA does not protect against physical device seizure where the TPM is
        bypassed at the hardware level.
  \item EIDA does not provide anonymity---\texttt{uPriv} is permanently linked
        to user identity.
  \item EIDA does not address CA availability---a downed CA blocks
        transactions.
\end{itemize}

% ═══════════════════════════════════════════════════════════════════════════
\section{Comparison with Existing Approaches}

\subsection{Classical PKI}

\begin{table}[ht]
\centering
\caption{EIDA vs.\ Classical PKI}
\label{tab:pki}
\begin{tabular}{lll}
  \toprule
  \textbf{Property}       & \textbf{Classical PKI} & \textbf{EIDA}      \\
  \midrule
  Identity guarantee      & Static                 & Per-transaction    \\
  Key reuse               & Yes                    & Never              \\
  \texttt{uPriv} exposure & High                   & None               \\
  Replay protection       & External               & Native             \\
  CA breach impact        & Critical               & Minimal            \\
  Forward secrecy         & Session-level          & Identity-level     \\
  Centralisation          & High                   & None               \\
  \bottomrule
\end{tabular}
\end{table}

\subsection{TLS Ephemeral Keys (TLS 1.3)}

TLS~1.3~\cite{RFC8446} introduces ephemeral session keys for forward
secrecy.  However, ephemeral keys apply to the \emph{session} layer, not the
identity layer.  The server certificate remains static and is reused across
all sessions; compromise of that certificate affects all past and future
sessions.  EIDA applies the ephemeral principle to the \emph{identity layer
itself}---the layer TLS never touches.

\subsection{One-Time Signature Schemes (Lamport, WOTS)}

One-time signature schemes share the single-use property with ePriPub.
However, they were designed for post-quantum resistance rather than identity
isolation, do not address the centralisation problem, and have no isolation
zone model or stateless CA mechanism.  ePriPub is architecturally distinct:
its purpose is identity-level forward secrecy, not quantum resistance.

\subsection{Decentralised Identity (DID, W3C)}

The W3C DID standard~\cite{W3CDID} moves identity off centralised registries.
However, DID does not address per-transaction key isolation, DID keys are
reused across transactions, and no single-use consent-token primitive exists
in the DID model.  EIDA is complementary to DID---ePriPub could serve as the
signing primitive within a DID framework.

% ═══════════════════════════════════════════════════════════════════════════
\section{Conclusion}

The fundamental vulnerability of current internet security architecture is
concentration.  When identity verification depends on centralised systems,
those systems become high-value targets.  A single successful attack exposes
millions.  This is not an implementation failure---it is an architectural one.

EIDA addresses this at the architectural level by introducing three properties
that classical PKI does not provide simultaneously:
\begin{itemize}
  \item \textbf{Identity-level forward secrecy} --- each transaction uses a
        unique, derived key pair that is destroyed after use.
  \item \textbf{Zero \texttt{uPriv} exposure} --- the permanent private key
        never leaves the hardware-backed isolation zone.
  \item \textbf{Distributed attack surface} --- there is no central repository
        of user identity material worth attacking.
\end{itemize}

The core primitive, ePriPub, achieves these properties by combining
well-established cryptographic standards---HKDF~\cite{RFC5869} and
Ed25519~\cite{RFC8032}---in a novel architectural pattern.  No new
cryptographic assumptions are required.

EIDA does not replace existing infrastructure.  It extends it.  CA systems,
PKI chains, and existing identity frameworks remain valid---EIDA adds a
per-transaction isolation layer on top of them.

The shift EIDA proposes is conceptual as much as technical: security should
not be something done \emph{to} users by centralised authorities.  It should
be something that lives \emph{with} the user, on the user's device, under the
user's control.

% ═══════════════════════════════════════════════════════════════════════════
\bibliographystyle{alpha}
\begin{thebibliography}{10}

\bibitem{RFC5869}
H.~Krawczyk and P.~Eronen.
\newblock {HMAC-based Extract-and-Expand Key Derivation Function (HKDF)}.
\newblock RFC~5869, Internet Engineering Task Force, May 2010.
\newblock \url{https://www.rfc-editor.org/rfc/rfc5869}

\bibitem{RFC8032}
S.~Josefsson and I.~Liusvaara.
\newblock {Edwards-Curve Digital Signature Algorithm (EdDSA)}.
\newblock RFC~8032, Internet Engineering Task Force, January 2017.
\newblock \url{https://www.rfc-editor.org/rfc/rfc8032}

\bibitem{Bernstein2012}
D.~J.~Bernstein, N.~Duif, T.~Lange, P.~Schwabe, and B.-Y.~Yang.
\newblock High-speed high-security signatures.
\newblock \emph{Journal of Cryptographic Engineering}, 2(2):77--89, 2012.

\bibitem{W3CDID}
{W3C}.
\newblock {Decentralized Identifiers (DIDs) v1.0}.
\newblock W3C Recommendation, July 2022.
\newblock \url{https://www.w3.org/TR/did-core/}

\bibitem{RFC8446}
E.~Rescorla.
\newblock {The Transport Layer Security (TLS) Protocol Version 1.3}.
\newblock RFC~8446, Internet Engineering Task Force, August 2018.
\newblock \url{https://www.rfc-editor.org/rfc/rfc8446}

\end{thebibliography}

\end{document}
